% ═══════════════════════════════════════════════════════════════
% LEHRERHINWEISE: Waagerechter Wurf (Physik Klasse 10E)
% Version 3.0 – Passend zu AB-01/LP-01 Version 3.0
% ═══════════════════════════════════════════════════════════════

\documentclass[10pt, a4paper]{article}

\usepackage[utf8]{inputenc}
\usepackage[T1]{fontenc}
\usepackage[ngerman]{babel}
\usepackage{helvet}
\renewcommand{\familydefault}{\sfdefault}

\usepackage[margin=1.5cm, top=1.8cm, bottom=1.5cm]{geometry}
\usepackage{enumitem}
\usepackage{tabularx}
\usepackage{booktabs}
\usepackage{array}
\usepackage{amsmath}
\usepackage{amssymb}

\usepackage{xcolor}
\usepackage{tcolorbox}
\usepackage{fancyhdr}
\usepackage{lastpage}

\setlist{nosep, leftmargin=1.5em}

% ═══════════════════════════════════════════════════════════════
% FARBEN
% ═══════════════════════════════════════════════════════════════

\definecolor{physik}{HTML}{2c5aa0}
\definecolor{grau}{HTML}{888888}
\definecolor{hellgrau}{HTML}{f5f5f5}
\definecolor{loesung}{HTML}{c62828}
\definecolor{hinweis}{HTML}{b35c00}

\newcommand{\fachfarbe}{physik}

% ═══════════════════════════════════════════════════════════════
% VARIABLEN
% ═══════════════════════════════════════════════════════════════

\newcommand{\thema}{Waagerechter Wurf}
\newcommand{\fach}{Physik}
\newcommand{\klasse}{10E}

% ═══════════════════════════════════════════════════════════════
% KOPF- UND FUSSZEILE
% ═══════════════════════════════════════════════════════════════

\pagestyle{fancy}
\fancyhf{}
\renewcommand{\headrulewidth}{0.4pt}
\renewcommand{\footrulewidth}{0pt}

\lhead{\small\fach{} \klasse{} -- \thema{} \textcolor{hinweis}{\textbf{[LEHRERHINWEISE]}}}
\rhead{\small DS 1 · 2. Halbjahr}

\cfoot{\small\thepage/\pageref{LastPage}}

% ═══════════════════════════════════════════════════════════════
% BEFEHLE
% ═══════════════════════════════════════════════════════════════

\newcommand{\abschnitt}[1]{%
    \vspace{8pt}
    \noindent\textbf{\textcolor{\fachfarbe}{#1}}
    \vspace{4pt}
}

\newtcolorbox{hinweisbox}[1][]{
    colback=orange!5,
    colframe=hinweis,
    boxrule=0.8pt,
    arc=0pt,
    left=6pt, right=6pt, top=6pt, bottom=6pt,
    fonttitle=\bfseries\small,
    title=#1
}

\newtcolorbox{loesungbox}{
    colback=red!5,
    colframe=loesung,
    boxrule=0.8pt,
    arc=0pt,
    left=6pt, right=6pt, top=6pt, bottom=6pt
}

% ═══════════════════════════════════════════════════════════════
% DOKUMENT
% ═══════════════════════════════════════════════════════════════

\begin{document}

\begin{center}
{\large\bfseries Lehrerhinweise: \thema}\\[2pt]
{\small\textcolor{grau}{Kurzplan + Lösungen + Differenzierung}}
\end{center}

\vspace{-2pt}
\hrule
\vspace{8pt}

% ═══════════════════════════════════════════════════════════════
% ÜBERBLICK
% ═══════════════════════════════════════════════════════════════

\abschnitt{Überblick}

\begin{tabular}{@{}ll@{}}
\textbf{Fach/Klasse:} & Physik 10E (Gymnasium) \\
\textbf{Thema:} & Waagerechter Wurf \\
\textbf{Position:} & DS 1 im 2. Halbjahr (nach Kinematik-Klausur) \\
\textbf{Zeit:} & 90 Minuten (Doppelstunde) \\
\textbf{Rahmenplan:} & Überlagerung von Bewegungen \\
\textbf{AB:} & 24 BE (K1--K5, Ü1--Ü4) \\
\end{tabular}

\vspace{8pt}

\abschnitt{Feinziele}

\begin{enumerate}[label=FZ\arabic*:]
    \item SuS beschreiben das Überlagerungsprinzip (AFB I)
    \item SuS berechnen Fallzeit aus Höhe mit $t = \sqrt{2h/g}$ (AFB II)
    \item SuS berechnen Wurfweite mit $s_x = v_0 \cdot t$ (AFB II)
    \item SuS berechnen Aufprallgeschwindigkeit $v_{\text{res}}$ und -winkel $\alpha$ (AFB II)
    \item SuS bestimmen $v_0$ durch Formelumstellung (AFB II)
    \item SuS beschreiben schrägen Wurf qualitativ als Überlagerung (AFB I)
    \item SuS erklären Orbitalmechanik als kontinuierlichen waagerechten Wurf (AFB III)
\end{enumerate}

% ═══════════════════════════════════════════════════════════════
% STUNDENVERLAUF
% ═══════════════════════════════════════════════════════════════

\abschnitt{Stundenverlauf (90 Min)}

\begin{tabular}{|c|p{2.2cm}|p{5.5cm}|c|c|}
\hline
\textbf{Min} & \textbf{Phase} & \textbf{Inhalt} & \textbf{SF} & \textbf{Material} \\
\hline
0--8 & Einstieg & Gedankenexperiment: Fallball vs. Wurfball & PL & LP 1--2 \\
\hline
8--15 & Beobachtung & SIM ausführen, Vergleichsball aktivieren & EA & LP 2, SIM Tab 1 \\
\hline
15--25 & Erarbeitung & Überlagerungsprinzip erklären & EA & LP 3--5 $\to$ AB K1 \\
\hline
25--35 & Formeln & $t$, $s_x$ einführen & EA & LP 5--6 $\to$ AB K2 \\
\hline
35--45 & Übung 1 & Grundaufgabe rechnen ($h=20$m, $v_0=15$m/s) & EA & LP 6--7 $\to$ AB Ü1 \\
\hline
45--55 & Aufprall & $v_y$, $v_{\text{res}}$, Challenge: tan-Formel & EA & LP 8--10 $\to$ AB K3, Ü2 \\
\hline
55--65 & Übung 2 & Anwendung Hilfspaket ($h=45$m) & EA/PA & LP 11 $\to$ AB Ü3 \\
\hline
65--72 & Transfer & Umkehraufgabe ($t$, $s_x$ gegeben) & EA & LP 12 $\to$ AB Ü4 \\
\hline
72--80 & Erweiterung & Schräger Wurf (qualitativ) & EA & LP 13 $\to$ AB K4 \\
\hline
80--87 & Ausblick & Orbitalmechanik, ISS-Erklärung & EA & LP 14 $\to$ AB K5 \\
\hline
87--90 & Sicherung & Selbstcheck, Zusammenfassung & PL & LP 15, AB \\
\hline
\end{tabular}

\vspace{8pt}

% ═══════════════════════════════════════════════════════════════
% HÄUFIGE FEHLER
% ═══════════════════════════════════════════════════════════════

\abschnitt{Häufige Schülerfehler}

\begin{hinweisbox}[Typische Denkfehler]
\begin{enumerate}
    \item \textbf{``Der geworfene Ball braucht länger, weil er weiter fliegt.''}

    $\to$ Horizontale Strecke $\neq$ Zeit. Die Fallzeit hängt nur von $h$ ab.

    \item \textbf{Wurzel vergessen:} SuS rechnen $t = 2h/g$ statt $t = \sqrt{2h/g}$.

    $\to$ An Herleitung aus $s = \frac{1}{2}gt^2$ erinnern.

    \item \textbf{Tan-Formel verwechselt:} $\tan(\alpha) = v_x/v_y$ statt $v_y/v_x$.

    $\to$ Challenge-Aufgabe im LP hilft, richtige Formel selbst herzuleiten.

    \item \textbf{Horizontale Geschwindigkeit ändert sich:} SuS denken, $v_x$ nimmt ab.

    $\to$ Keine horizontale Kraft $\Rightarrow$ keine horizontale Beschleunigung.

    \item \textbf{Pythagoras-Fehler:} $v_{\text{res}} = v_x + v_y$ statt $\sqrt{v_x^2 + v_y^2}$.

    $\to$ Tripel-Tabelle auf AB nutzen (3-4-5, 15-20-25 etc.).
\end{enumerate}
\end{hinweisbox}

% ═══════════════════════════════════════════════════════════════
% DIFFERENZIERUNG
% ═══════════════════════════════════════════════════════════════

\abschnitt{Differenzierung}

\begin{tabular}{|l|l|p{6.5cm}|}
\hline
\textbf{Niveau} & \textbf{Aufgaben} & \textbf{Charakteristik} \\
\hline
Basis & Ü1, Ü2 a/b & Direktes Einsetzen, bekannte Formeln \\
\hline
Standard & Ü2 c, Ü3 & Mehrere Schritte, Aufprallwinkel bestimmen \\
\hline
Erweitert & Ü4, K4, K5 & Formelumstellung, Transfer, qualitativ \\
\hline
\end{tabular}

\vspace{4pt}
\textbf{Empfehlung:} Alle SuS bearbeiten K1--K3 und Ü1--Ü2. Schnelle SuS zusätzlich Ü3--Ü4 und K4--K5.

\vspace{8pt}

% ═══════════════════════════════════════════════════════════════
% LÖSUNGEN
% ═══════════════════════════════════════════════════════════════

\abschnitt{Kurzlösungen AB-01}

\begin{loesungbox}
\small
\textbf{K1 (Überlagerungsprinzip):}
\begin{itemize}[nosep]
    \item Horizontal: \textbf{gleichförmige Bewegung}
    \item Vertikal: \textbf{freier Fall}
    \item beeinflusst die \textbf{Fallzeit} nicht
    \item landen \textbf{gleichzeitig}
\end{itemize}

\vspace{4pt}
\textbf{K2 (Formeln):}
\begin{itemize}[nosep]
    \item $t = \sqrt{2h/g}$
    \item $s_x = v_0 \cdot t$
\end{itemize}

\vspace{4pt}
\textbf{K3 (Aufprall):}
\begin{itemize}[nosep]
    \item $v_y = g \cdot t$
    \item $v_{\text{res}} = \sqrt{v_x^2 + v_y^2}$
    \item $\tan(\alpha) = v_y / v_x$
\end{itemize}

\vspace{4pt}
\textbf{K4 (Schräger Wurf):}
\begin{itemize}[nosep]
    \item 1. gleichförmige Horizontalbewegung
    \item 2. senkrechter Wurf nach oben
    \item Steigzeit \textbf{=} Fallzeit
    \item Max. Wurfweite bei \textbf{45°}
    \item Bahnform: \textbf{Parabel}
\end{itemize}

\vspace{4pt}
\textbf{K5 (Orbitalmechanik):}
\begin{itemize}[nosep]
    \item Orbit = kontinuierlicher \textbf{waagerechter Wurf}
    \item $v_1 \approx$ \textbf{7,9} km/s (Orbitgeschwindigkeit)
    \item $v_2 \approx$ \textbf{11,2} km/s (Fluchtgeschwindigkeit)
    \item ISS: fällt ständig, aber Erde krümmt sich weg
\end{itemize}
\end{loesungbox}

\vspace{4pt}

\begin{loesungbox}
\small
\textbf{Ü1:} $h = 20\,\text{m}$, $v_0 = 15\,\text{m/s}$ \hfill (4 BE)
\begin{itemize}[nosep]
    \item a) $t = \sqrt{40/10} = \mathbf{2\,\text{s}}$
    \item b) $s_x = 15 \cdot 2 = \mathbf{30\,\text{m}}$
\end{itemize}

\vspace{4pt}
\textbf{Ü2:} (gleiche Werte wie Ü1) \hfill (6 BE)
\begin{itemize}[nosep]
    \item a) $v_y = 10 \cdot 2 = \mathbf{20\,\text{m/s}}$
    \item b) $v_{\text{res}} = \sqrt{15^2 + 20^2} = \sqrt{625} = \mathbf{25\,\text{m/s}}$ \quad (3-4-5 $\times 5$!)
    \item c) $\tan(\alpha) = 20/15 = 1{,}33 \Rightarrow \alpha \approx \mathbf{53°}$
\end{itemize}

\vspace{4pt}
\textbf{Ü3:} $h = 45\,\text{m}$, $v_0 = 20\,\text{m/s}$ \hfill (10 BE)
\begin{itemize}[nosep]
    \item a) $t = \sqrt{90/10} = \sqrt{9} = \mathbf{3\,\text{s}}$
    \item b) $s_x = 20 \cdot 3 = \mathbf{60\,\text{m}}$
    \item c) $v_y = 10 \cdot 3 = \mathbf{30\,\text{m/s}}$
    \item d) $v_{\text{res}} = \sqrt{400 + 900} = \sqrt{1300} \approx \mathbf{36\,\text{m/s}}$
    \item e) $\tan(\alpha) = 30/20 = 1{,}5 \Rightarrow \alpha \approx \mathbf{56°}$ (zwischen 53° und 60°)
\end{itemize}

\vspace{4pt}
\textbf{Ü4:} $t = 4\,\text{s}$, $s_x = 100\,\text{m}$ \hfill (4 BE)
\begin{itemize}[nosep]
    \item a) $v_0 = s_x / t = 100/4 = \mathbf{25\,\text{m/s}}$
    \item b) $h = \frac{1}{2} g t^2 = \frac{1}{2} \cdot 10 \cdot 16 = \mathbf{80\,\text{m}}$
\end{itemize}
\end{loesungbox}

\vspace{8pt}

% ═══════════════════════════════════════════════════════════════
% MATERIAL
% ═══════════════════════════════════════════════════════════════

\abschnitt{Benötigtes Material}

\begin{itemize}[nosep]
    \item Tablets/Laptops (1 pro SuS oder 1 pro 2 SuS)
    \item Arbeitsblatt AB-01 (ausgedruckt, 1 pro SuS)
    \item Lernpfad LP-01 (digital)
    \item Simulation SIM-01 (im LP eingebettet, Tabs 1--3)
    \item Optional: Zwei gleichschwere Bälle für Live-Demo
\end{itemize}

\vspace{8pt}

\abschnitt{Vorbereitung}

\begin{enumerate}[nosep]
    \item WLAN funktionsfähig?
    \item LP-URL an die Tafel schreiben oder QR-Code bereitstellen
    \item AB für alle kopieren (2 Seiten, doppelseitig möglich)
    \item Optional: Live-Experiment vorbereiten
\end{enumerate}

\vspace{8pt}

\abschnitt{Hinweis zur Challenge (LP Schritt 9)}

Die Challenge-Aufgabe im Lernpfad lässt SuS die Tangens-Formel selbst aus 10 ähnlichen Optionen auswählen. \textbf{Keine Versuchsbegrenzung}, aber auch keine Hinweise. Dies fördert eigenständiges Denken und Transferleistung.

\vspace{8pt}

\abschnitt{Stundenleistung MT-01 (optional)}

Nach der DS kann die Stundenleistung MT-01 (27 BE, 20 Min) als digitaler Test durchgeführt werden. Alle Aufgaben sind auto-auswertbar und kopfrechenbar. \textbf{Hilfsmittel:} AB + LP (open book).

\end{document}
